\subsection{Data and graph representation}
\label{sec:dados_grafos}

The vertices in the file names were not in the sensor frame, and the \num{1152} end-cap channels had flipped signs in $y$ and $z$; both corrections were applied during the construction of the dataset and checked on the \num{9364} events. Three hot channels, which had counts in every event and held \qty{47.1}{\percent} of the raw counts, were replaced by the mean of their healthy neighbours: the dataset stores the corrected light. One position was simulated twice, and three files were copies of others. After the corrections, the centroid of the detected light follows the vertex, with a Pearson correlation coefficient $r$ of \num{0.908} in $x$, \num{0.871} in $y$ and \num{0.999} in $z$ over the \num{9364} events.

The events were split into \num{7490} training, \num{937} validation and \num{937} test events, stratified in \num{9} bins per axis (\num{729} groups), a split shared by all models. Given the symmetry of the channels under sign flips of the axes, each training event was mirrored into the other octants: \num{56884} training rows, \num{7.59} per event; the validation and test events stay in the simulated octant. The test vertices lie, on average, \qty{49.31}{\cm} from the nearest training vertex and, except for the position simulated twice, are grid positions absent from the training set: the errors measure interpolation on a \qty{0.5}{\m} grid, with one photon draw per event.

The nodes of the graph $G = (\mathcal{V}, \mathcal{E})$ are the $|A|$ active channels of the event, \num{245} on average in the test events, and one global node \citep{gilmer2017neural} at the light centroid. Each channel receives edges from the $k = {}$\num{4} nearest active channels, in a $k$-nearest-neighbour (\mbox{$k$-NN}) graph, and exchanges one edge in each direction with the global node: six directed edges per channel (\num{248} nodes and \num{1482} edges in the event of Figure~\ref{fig:evento_exemplo}). Each node $i$ has \num{17} features $\mathbf{f}_i$: position, light, light fraction, distance and direction to the centroid, direction to the brightest channel, light of the neighbours, eigenvalues of the light-weighted covariance matrix of the positions (common to all nodes of the event) and light rank. Each edge has \num{23} features $\mathbf{e}_{ji}$: distance, direction, inverse square of the distance, light difference and light similarity, and \num{16} radial basis functions of the distance up to \qty{7}{\m}. The \num{33} event-level scalars, $\mathbf{u}$, are centroids with linear and quadratic weights, total light, number of active channels, largest channel light, spreads, skewnesses and position quantiles. The normalisations use only the training set (\num{8192} sampled graphs), and each vertex coordinate is standardised.
